% resume_content_v2.tex -- Updated March 2026
% Shared body file, included via \input{resume_content_v2} from main .tex files.
% Requires \newtoggle{detailedVersion} to be defined before \input.

\section{\centerline{SUMMARY}}
\vspace{12pt}
\center{\it{Hands-on engineering leader with 20 years shipping AI, robotics, and simulation from prototype to production}}

\vspace{0.2in} 
\section{\centerline{ CORE EXPERTISE }}
\vspace{8pt}
\begin{table}[ht]
\centering
\begin{tabular}{lll}
\textbullet~Artificial Intelligence &  \textbullet~Robotics \& Control Systems &  \textbullet~Machine Learning \\
\textbullet~Deep Learning & \textbullet~Generative AI / Foundation Models & \textbullet~Computer Vision \\
\textbullet~Autonomous Vehicles / ADAS & \textbullet~Simulation \& HIL/SIL/MIL & \textbullet~Planning under Uncertainty \\
\textbullet~Real-Time \& GPU Systems & \textbullet~Embedded Systems & \textbullet~Signal \& Image Processing \\
\textbullet~Systems Software & \textbullet~Building \& Scaling Teams & \textbullet~Engineering Leadership
\end{tabular}
\end{table}
 
% ============================================================
% WORK EXPERIENCE
% ============================================================
\vspace{0.2in} 
\section{\centerline{WORK EXPERIENCE}}

\vspace{8pt}
{\sl \href{https://wayve.ai}{\textbf{Wayve}}} \hfill        August 2025 - Present \\
\href{https://wayve.ai}{\textbf{San Francisco Bay Area}}       \hfill   \textbf{Technical Lead, Senior Staff Engineer}
   \begin{itemize} \itemsep -2pt
   \item Building Wayve's cutting-edge closed-loop simulation based on generative AI foundation models
   \begin{itemize}
   {\item[$\checkmark$] Helped develop Wayve's simulation strategy}
   {\item[$\checkmark$] Collaborated on Verification and Release approaches (e.g.\ HIL validation, PIL)}
   {\item[$\checkmark$] Developed and helped maintain core simulation engines for open-loop and closed-loop simulation}
   {\item[$\checkmark$] Implemented simulation support for AI evaluation of safety and L2+ product features}
   \end{itemize}
\end{itemize}

\vspace{8pt}
{\sl \href{http://applied.co}{\textbf{Applied Intuition}}} \hfill        January 2020 - July 2025 \\
\href{http://applied.co}{\textbf{Mountain View, California}}       \hfill   \textbf{Engineering Manager}
   \begin{itemize} \itemsep -2pt
   \item Founded, created and led team building X-In-the-Loop (XIL) simulation products at Applied Intuition
   \begin{itemize}
   {\item[$\checkmark$] Founding XIL member; led Object HIL from concept to production in under one year}
	{\item[$\checkmark$] Managed distributed engineering across Mountain View, Austin, and Munich}
	{\item[$\checkmark$] Built team with depth in C++, real-time systems, GPU programming, and simulation}
   {\item[$\checkmark$] Owned product vision and technical roadmap since the team was founded}
   {\item[$\checkmark$] Shipped Object HIL, Sensor HIL, and Replay HIL}
   {\item[$\checkmark$] Built the Model-in-the-Loop (MIL) toolbox for the Applied Development Platform (ADP)}
   {\item[$\checkmark$] Designed virtual ECU (vECU) simulator for software-in-the-loop (SIL) testing}
   {\item[$\checkmark$] Contributed to core simulator engine and AD/ADAS validation interfaces}
   \end{itemize}
   \item Responsible for overall architecture, product direction and engineering design
   \begin{itemize}
   {\item[$\checkmark$] Partnered broadly across Applied engineering---XIL touches most of the org}
   {\item[$\checkmark$] Helped hire Applied's vehicle platform team from scratch}
   {\item[$\checkmark$] Interviewed senior robotics, vehicle hardware, controls, and planning candidates (beyond GenSWE)}
   {\item[$\checkmark$] Player/coach: primary architect for XIL; reviews and hands-on coding as needed}
   \end{itemize}
\end{itemize}


\vspace{8pt}
{\sl \href{http://artificial.com}{\textbf{Artificial, Inc.}}} \hfill        September 2019 - November 2019 \\
\href{http://artificial.com}{\textbf{Palo Alto, California}}       \hfill   \textbf{Consultant}
   \begin{itemize} \itemsep -2pt
   \item Helped refine product-market fit for Robotic Lab Automation
   \begin{itemize}
   {\item[$\checkmark$] Partnered with product and CTO on roadmap, features, and architecture}
	{\item[$\checkmark$] Built three.js simulator prototype to visualize robotic lab workcells}
   {\item[$\checkmark$] Prototype became a core product surface---customers preview how assays will be automated}
   \end{itemize}
\end{itemize}

\vspace{8pt}
{\sl \href{https://web.archive.org/web/20230130130556/https://level.ai/}{\textbf{LEVEL.ai}}} \hfill        August 2016 - August 2019 \\
\href{https://web.archive.org/web/20230130130556/https://level.ai/}{\textbf{Palo Alto, California}}       \hfill   \textbf{Founder, CTO}
   \begin{itemize} \itemsep -2pt
   \item Co-founded company in my garage, with Leonard Speiser and Nick Leindecker
   \begin{itemize}
   	{\item[$\checkmark$] We raised 25 million in 2 rounds (Series A and B)}
	{\item[$\checkmark$] Board member and part of management team making planning/prioritization/hiring/business decisions}
	{\item[$\checkmark$] Curator-in-chief of the company's culture and core values}
   \end{itemize}
   \item Responsible for overall architecture, R\&D and software
   \begin{itemize}
   {\item[$\checkmark$] Named inventor on 5 issued US patents covering AI-driven cooking, energy steering, and splatter prevention (see Patents section)}
   {\item[$\checkmark$] Developed core AI: Planning \& control algorithms and led development of deep learning perception stack}
   {\item[$\checkmark$] Helped build our engineering team for machine learning, computer vision, sensing and planning}
   {\item[$\checkmark$] Learned to tradeoff between building for rapid iteration and operational stability/scalability}
   {\item[$\checkmark$] Player/coach - managed engineers, was chief architect, did code-reviews, while being hands-on throughout}
   {\item[$\checkmark$] Our stack involves hardware, systems-software, Android applications, Deep learning, data-pipelines and cloud apps.}
   \end{itemize}
\end{itemize}


\vspace{8pt}
{\sl \href{http://www.clover.com}{\textbf{Clover Network Inc.}}} \hfill        October 2013 - September 2016 \\
\href{http://www.clover.com/team}{\textbf{Mountain View, California}}       \hfill   \textbf{Software Engineer, Data Scientist}
   \begin{itemize} \itemsep -2pt
   \item Built and led Clover's internal data science team. Some projects we worked on were:
	  \begin{itemize}
		  {\item[$\checkmark$] Merchant retention modeling}
		  {\item[$\checkmark$] Real time and forensic failure analytics}
		  {\item[$\checkmark$] A/B testing (identification and design of experiments)}
		  {\item[$\checkmark$] App market suggestions}
	  \end{itemize}	  
   \item High-performance \href{http://www.emvco.com/}{EMVCo contactless NFC} through \href{https://www.clover.com/pos-hardware/mobile}{LCD display}
   \begin{itemize}
   \item[$\checkmark$]
   Played pivotal role in Clover achieving analog and digital L1-type approvals
   \iftoggle{detailedVersion}{\item[$\checkmark$] NFC (reader) through display required fighing Physics - we were the world's first to overcome challenge
   \begin{itemize}\item[\tiny$\blacksquare$] Led cross-functional team (Hardware, Systems-software, RF-experts) realize what seemed like mission impossible
   \item[\tiny$\blacksquare$] Demonstrated ability to jump into unknown domains, work with cross-functional experts, and rapidly deliver on a hard project
   \end{itemize}}{}
   \end{itemize}

   \item Systems Software and Firmware Engineering
   \begin{itemize}
   \item[$\checkmark$] Potential to work at any level of systems stack
   \iftoggle{detailedVersion}{\item[$\checkmark$] Developed portions of secure payments software
   \begin{itemize}\item[\tiny$\blacksquare$] Developed Linux kernel drivers for custom hardware interfaces
   \item[\tiny$\blacksquare$] Worked on Android OS framework, Android recovery etc.
   \item[\tiny$\blacksquare$] Code written mostly in C, Java, C++ and scripting languages such as Python and bash
   \end{itemize}}{}
   \end{itemize}

   \item Developed barcode scanners for Clover Point of Sales terminals \itemsep -2pt
   \begin{itemize}
   \item[$\checkmark$]  \itemsep -8pt%
    Scanner enables rapid product barcode and QR-code scanning using low-cost camera
    \iftoggle{detailedVersion}{\begin{itemize}\item[\tiny$\blacksquare$] Combines techniques from Computer Vision, Image Processing and Algorithms
   \item[\tiny$\blacksquare$] All code written in C++, Java for Android, with prototyping in Matlab
   \item[\tiny$\blacksquare$] Used hardware parallelization, OpenCV and developed algorithms for improved image enhancement
   \end{itemize}}{}
     \end{itemize}
  \item Implemented features in various core Clover software  \itemsep -2pt
   \begin{itemize}
   \item[$\checkmark$]  \itemsep -8pt%
    Added features to Clover's server, Android platform, Android apps and web dashboard
    \iftoggle{detailedVersion}
    { \begin{itemize}
       \item[\tiny$\blacksquare$] Languages used are Java, C++, C, Javascript and SQL
      \end{itemize}
    } {}
   \end{itemize}
   \item Named inventor on 3 issued US patents covering NFC filter tuning and accessible PIN entry (see Patents section)
\end{itemize}
 
\vspace{8pt}
{\sl \href{http://www.usc.edu}{\textbf{University of Southern California, Los Angeles}}} \hfill        January 2006 - September 2013 \\
\href{http://robotics.usc.edu/resl}{\textbf{Robotic Embedded Systems Laboratory}}       \hfill   \textbf{Research Assistant}
   \begin{itemize} \itemsep -2pt
   \item \href{http://robotics.usc.edu/~ampereir/wordpress/?page_id=131}{Developed  path planners for glider AUVs under High Uncertainty} \itemsep -2pt
   \begin{itemize}\item[$\checkmark$]  \itemsep -8pt%
    1st planners for glider AUVs to deal with uncertainty in ocean model predictions 
    \iftoggle{detailedVersion}
    {\begin{itemize}\item[\tiny$\blacksquare$] Combines techniques from AI, Machine Learning and Robotics
   \item[\tiny$\blacksquare$] All code written in C++ and Python (STL, Boost, NumPy/SciPy)
   \item[\tiny$\blacksquare$] Used hardware parallelization and advanced data structures to speed up computations
    \end{itemize}
   } {}
   \end{itemize}
    
   \item \href{http://robotics.usc.edu/~ampereir/wordpress/?page_id=134}{Developed a long-range RF communication network for AUVs in the Southern California region}
   \begin{itemize} \item[$\checkmark$]Improved data throughput by over 24x and reduced operation cost by over 50\%
   \iftoggle{detailedVersion}{\begin{itemize}
   \item[\tiny$\blacksquare$]  Required Computer Networking, Embedded Systems and Signal Processing
   \item[\tiny$\blacksquare$]  All code written in C++ (both ARM and PC-based Linux) 
   \end{itemize}}{}
   \end{itemize}
   
    \item Developed a novel TCP-like communication protocol which adapts dynamically to sea state
   \begin{itemize} \item[$\checkmark$]1st motion-aware adaptive communication protocol
   \iftoggle{detailedVersion}{\begin{itemize}
   \item[\tiny$\blacksquare$]  Required Computer Networking, Reinforcement Learning and Signal Processing
   \item[\tiny$\blacksquare$]  All code written in C++ (both ARM and PC-based Linux) 
   \end{itemize}}{}
   \end{itemize}
   
    \item \href{http://robotics.usc.edu/~ampereir/wordpress/?page_id=365}{Developed the Guidance, Navigation and Control for Autonomous Surface Vehicles  including stereo-vision based obstacle avoidance} and \href{http://robotics.usc.edu/~ampereir/wordpress/?page_id=373}{Station-Keeping}
 \iftoggle{detailedVersion}{\begin{itemize} \item[$\checkmark$] Took only 2 months to complete autonomy from RC-control to Mission control}{}
    \iftoggle{detailedVersion}{
   \begin{itemize}%
   \item[\tiny$\blacksquare$]  Required Control Systems, Robotics, Computer Networking, Computer Vision, Image Processing
   \iftoggle{detailedVersion}{
   \item[\tiny$\blacksquare$]  All code on robot written in C++ (Linux); OpenCV for vision. GUI written in C\# .Net (Windows); Prototyping in Matlab (Control Systems Toolbox, System Identification Toolbox) and Simulink}{}
    \end{itemize}
    }{}
    \iftoggle{detailedVersion}{\end{itemize}}{}
 \end{itemize}

\vspace{8pt}
{\sl \href{http://welcome.isr.ist.utl.pt/home/}{\textbf{Institute for Systems and Robotics, Lisbon, Portugal}}} \hfill May 2008--August 2008  \\[2pt]
\href{http://www.grex-project.eu/}{\textbf{EU-GREX Project}} \hfill   \textbf{ Lead Developer}
\begin{itemize} \itemsep -2pt
 \item \href{http://robotics.usc.edu/~ampereir/wordpress/?page_id=397}{Implemented GREX Interface module for EU-GREX cooperative marine vehicle project}
   \begin{itemize}
   \item[$\checkmark$] Met all objectives on time despite large pre-existing code-base, and multi-national teams
   \item[$\checkmark$] Only team with a working robot at the Azores field test
   \end{itemize}
\end{itemize} 

\vspace{8pt}
{\sl \href{http://nio.org}{\textbf{National Institute of Oceanography, Goa, India}}}       \hfill               June 2002 - July 2005 \\[2pt]
\textbf{Marine Instrumentation \& Computer Division} \hfill \textbf{Research Fellow / Project Assistant}
 \begin{itemize} \itemsep -2pt
 \item  \href{http://robotics.usc.edu/~ampereir/wordpress/?page_id=425}{Led software development and hardware design on India's first Autonomous Underwater Vehicle (MAYA)}
 \begin{itemize} 
 \item[$\checkmark$] Wrote over 95\% of Control, Communications, Recovery, Guidance and Navigation code
 \item[$\checkmark$] Designed embedded circuitry for CAN actuator nodes
 \end{itemize}
 \item  Developed oceanographic instruments (wave-recorder with GSM telemetry, tide-gauges for UNESCO/IOC)
 \begin{itemize}
 \item[$\checkmark$] Tide gauges remained functional for over a decade without servicing
 \end{itemize}
 \item  Designed heading controller for ROSS autonomous surface craft; tested at sea
 \end{itemize}

% ============================================================
% PATENTS (verified via Google Patents, March 2026)
% ============================================================
\vspace{0.2in}
\section{\centerline{PATENTS}}
\vspace{8pt}
{\bf Markov Corporation / Level.ai} (5 issued US patents, 3 patent families)
\begin{itemize} \itemsep -2pt
\item US 10,009,957 B2 -- Electronic Oven with Infrared Evaluative Control (Jun 2018)
\item US 10,681,776 B2 -- Electronic Oven with Infrared Evaluative Control, continuation (Jun 2020)
\item US 11,632,826 B2 -- Electronic Oven with Infrared Evaluative Control, continuation (Apr 2023)
\item US 10,219,330 B2 -- Electronic Oven with Splatter Prevention (Feb 2019)
\item US 10,980,088 B2 -- Energy Absorption Monitoring for Intelligent Electronic Oven (Apr 2021)
\end{itemize}
{\bf Clover Network Inc.} (3 issued US patents, 2 patent families)
\begin{itemize} \itemsep -2pt
\item US 9,319,029 B1 -- System and Method for Automatic Filter Tuning (Apr 2016)
\item US 9,716,482 B2 -- System and Method for Automatic Filter Tuning, continuation (Jul 2017)
\item US 9,483,693 B1 -- Free-Hand Character Recognition on a Touch Screen POS Terminal (Nov 2016)
\end{itemize}

% ============================================================
% ENTREPRENEURIAL EXPERIENCE
% ============================================================
\vspace{0.2in}
\section{\centerline{ENTREPRENEURIAL EXPERIENCE}}
\vspace{10pt}
\begin{itemize} \itemsep -2pt
\item Markov / Level.ai -- Technical co-founder of Markov Corporation (Level.ai), where I led engineering and R\&D.
\item \href{https://ampy.us}{ampy.us} -- Founded consulting practice (Aug 2019--Present) providing strategic and technical consulting for robotics, automation, machine learning, computer vision, and embedded systems.
\end{itemize}

% ============================================================
% MENTORING EXPERIENCE
% ============================================================
\section{\centerline{MENTORING EXPERIENCE}}
\vspace{10pt}
% TODO: Arvind to provide academic mentoring details here
% (e.g. interns, PhD students, university collaborators -- institutions, counts, outcomes)
\centerline{\it Placeholder -- to be updated with recent academic mentoring experience.}

% ============================================================
% CERTIFICATIONS
% ============================================================
\vspace{0.2in}
\section{\centerline{CERTIFICATIONS}}
\vspace{8pt}
\begin{itemize} \itemsep -2pt
\item Business Leadership -- Harrison Metal (Jan 2019)
\item General Management -- Harrison Metal (Nov 2018)
\end{itemize}

% ============================================================
% EDUCATION
% ============================================================
\vspace{0.2in}
\section{\centerline{EDUCATION}}
\vspace{8pt}
{\sl PhD}, Computer Science \\
{\bf University of Southern California, CA \hspace{0.5in}   \hfill Sep 2013 (graduated)}\\
Thesis - \href{http://robotics.usc.edu/~ampereir/wordpress/?page_id=131}{Path Planning for Autonomous Underwater Vehicles in the Open Ocean}\\
Advisor - \href{http://robotics.usc.edu/~gaurav}{Prof. Gaurav S. Sukhatme}

{\sl Master of Science}, Electrical Engineering \\
{\bf University of Southern California, CA \hspace{0.5in}   \hfill May 2007 (graduated)}\\
Thesis - \href{http://robotics.usc.edu/~ampereir/wordpress/?page_id=365}{Navigation and Guidance of an Autonomous Surface Vehicle} \\
Advanced DSP project - \href{http://robotics.usc.edu/~ampereir/wordpress/?p=156}{Tracking groups of people in real-time video using TI DSP processor} \\
Sensing and Planning in Robotics Project -  \href{http://robotics.usc.edu/~ampereir/pubs/report547.pdf}{Mapless Localization from Fused Sensor Data}
 
 \vspace{0.2in}
 
{\sl Bachelor of Engineering}, Electronics and Communications \\
{\bf \href{http://www.vtu.ac.in/en}{VTU}, Belgaum, Karnataka, India   \hspace{0.5in} 1st class with Distinction   \hfill    Jul 2002 (graduated)} \\
Project 1 - \href{http://robotics.usc.edu/~ampereir/wordpress/?p=382}{Development of a Automatic Weather Station} \\
Project 2 - \href{http://robotics.usc.edu/~ampereir/wordpress/?p=382}{PC-based Blood Sample Colorimeter}

% ============================================================
% LANGUAGES
% ============================================================
\vspace{0.2in}
\section{\centerline{LANGUAGES}}
\vspace{8pt}
\begin{table}[ht]
\centering
\begin{tabular}{ll}
\textbullet~English (native/bilingual) & \textbullet~Hindi (native/bilingual) \\
\textbullet~Konkani (native/bilingual) & \textbullet~Portuguese (limited working proficiency) \\
\textbullet~Marathi (limited working proficiency) &
\end{tabular}
\end{table}

\vspace{0in} 
